\section{Data}

\subsection*{Zdroj}

Projekt vychází z datasetu Steam Games Dataset~\cite{fronkon_dataset}, který obsahuje metadata
her získaná přes API Steamu. Použit je soubor \texttt{games.json} v pevné revizi, aby byla data
reprodukovatelná. Soubor obsahuje 148\,909 záznamů a 42 sloupců. Kromě metadat obsahuje
i odkazy na obrázky hry: header a seznam screenshotů.

\subsection*{Čištění}

Hry s nulovým počtem recenzí nelze použít, protože nulu nelze odlišit od chybějících dat:
3\,467 her má nulový počet recenzí, ačkoli má nenulový počet doporučení
(\texttt{recommendations}). Podíl takových chybějících hodnot navíc roste u novějších her.
Medián poměru počtu recenzí a počtu doporučení v letech 2011--2023 pozvolna klesá z 1,43 na 1,06,
od roku 2024 ale výrazně klesá (0,99 v roce 2024, 0,81 v roce 2025). Počty recenzí novějších her
jsou tedy podhodnocené, a proto se data omezují na hry vydané do roku 2023. Postup čištění
a počty her po jednotlivých krocích shrnuje tabulka~\ref{tab:cisteni}.

\begin{table}[h]
    \centering
    \begin{tabular}{lr}
        \toprule
        Krok & Počet her \\
        \midrule
        Původní data & 148\,909 \\
        Vydáno do roku 2023 & 77\,575 \\
        Alespoň 1 recenze & 64\,675 \\
        Alespoň 3 screenshoty a header & 64\,554 \\
        Bez softwaru (pouze softwarové žánry) & 63\,889 \\
        Bez duplicit (stejný název a vývojář) & 63\,804 \\
        \bottomrule
    \end{tabular}
    \caption{Počet her po jednotlivých krocích čištění.}
    \label{tab:cisteni}
\end{table}

U stejné hry vedené pod více identifikátory zůstává záznam s nejvíce recenzemi. Seznam
podporovaných jazyků obsahoval chybně zapsané položky (zbytky HTML a BBCode, více jazyků
v jedné položce), názvy jazyků se z nich proto extrahují regulárním výrazem.

\subsection*{Rozdělení dat}

Data jsou rozdělena časově podle roku vydání (tabulka~\ref{tab:rozdeleni}). Model se tak
vyhodnocuje na hrách vydaných později než trénovací data, což odpovídá praktickému použití
(predikce pro nově vydávané hry). Jednotkou rozdělení je hra, každá hra je tedy právě v jedné
množině.

\begin{table}[h]
    \centering
    \begin{tabular}{lcrr}
        \toprule
        Množina & Roky vydání & Počet her & Podíl \\
        \midrule
        Trénovací & do 2021 & 45\,565 & 71,4\,\% \\
        Validační & 2022 & 8\,611 & 13,5\,\% \\
        Testovací & 2023 & 9\,628 & 15,1\,\% \\
        \bottomrule
    \end{tabular}
    \caption{Časové rozdělení dat.}
    \label{tab:rozdeleni}
\end{table}

\subsection*{Vstupy modelu}

Metadata jsou zakódována do 77 vstupů:
\begin{itemize}
    \item standardizovaný rok vydání, $\log(1 + \text{cena})$ a $\log(1 + \text{počet jazyků})$,
    \item binární příznaky: hra zdarma, podpora Windows, macOS a Linuxu,
    \item multi-hot 22 žánrů s alespoň 50 hrami v trénovací množině,
    \item multi-hot 19 kategorií popisujících design hry (např. Single-player, Co-op, VR),
    \item multi-hot 29 jazyků podporovaných alespoň 1\,\% her v trénovací množině.
\end{itemize}

Obrazovým vstupem je header a první tři screenshoty každé hry. Obrázky jsou zmenšeny,
oříznuty na střed na čtverec $128 \times 128$ pixelů a jejich hodnoty škálovány do intervalu
$[0, 1]$.

\subsection*{Prevence úniku informací}

Vstupy obsahují pouze informace známé v době vydání hry. Vyřazeny jsou počty doporučení,
vlastníků, souběžných hráčů, herní doba, uživatelské tagy a počet DLC, které vznikají až po
vydání a souvisí s popularitou. Z kategorií jsou použity jen ty, které popisují design hry.
Kategorie přidávané až později (např. Steam Trading Cards nebo Remote Play) by popularitu
prozrazovaly. Ze stejného důvodu je vyřazen žánr Early Access, který odpovídá stavu hry
v době stažení dat. Parametry kódování (seznamy žánrů a jazyků, průměry a směrodatné
odchylky) se počítají pouze z trénovací množiny.

\subsection*{Omezení dat}

\begin{itemize}
    \item Počet recenzí je jen nepřímou mírou popularity. Starší hry měly na sběr recenzí
    více času, průměrné $y$ proto s rokem vydání klesá (v trénovací množině z 7,35 v roce 2013
    na 3,39 v roce 2021).
    \item Vyřazením her bez recenzí se model učí jen na hrách s alespoň jednou recenzí
    a pro hry bez recenzí jeho predikce neplatí.
    \item Rok vydání validačních a testovacích her leží mimo rozsah trénovací množiny
    (po standardizaci průměrně 1,51 a 1,90 směrodatné odchylky), model tedy musí extrapolovat.
    \item Obrázky odpovídají stavu stránky hry v době stažení. Vývojáři je mohli po vydání
    změnit, takže nemusí přesně odpovídat prezentaci v době vydání.
\end{itemize}